\clearpage
\begin{explanationbox}
\textbf{从初中底高到共底面积比：两个横坐标为什么相乘？}
\begin{center}% C088 PDF teaching figure; exact coordinates, equal x/y units.
\begin{tikzpicture}[x=1cm,y=1cm,>=stealth,font=\small]
\draw[rounded corners=5pt,black!25] (-1.2,-5.05) rectangle (11.1,4.3);
\node[font=\bfseries] at (4.95,3.7) {两次相似：横坐标比 $\to$ 斜边比 $\to$ 高之比};
\node at (1.5,2.7) {① 向 $x$ 轴投影，比较 $OB,OC$};
\node at (7.75,2.7) {② 共底 $OA$，比较两条高};
\begin{scope}
\coordinate (O) at (0,0); \coordinate (A) at (3,1.8); \coordinate (C) at (3,-1.8);
\coordinate (B) at (1.5,-.9); \coordinate (F) at (1.5,0); \coordinate (E) at (3,0);
\fill[orange!15] (O)--(A)--(C)--cycle;
\fill[teal!23] (O)--(B)--(F)--cycle;
\draw[->,black!45] (-.45,0)--(3.95,0) node[right] {$x$};
\draw[blue!65!black,thick] (O)--(A);
\draw[violet!75!black,dashed,thick] (O)--(C);
\draw[orange!85!black,very thick] (A)--(C);
\node[right=6pt] at (3,.65) {$2kx_1$};
\draw[teal!80!black,thick,dashed] (B)--(F);
\draw (1.32,0)--(1.32,-.18)--(1.5,-.18);
\draw (2.82,0)--(2.82,-.18)--(3,-.18);
\draw[black!65] (.55,0) arc[start angle=0,end angle=-30.96375653,radius=.55];
\fill (F) circle (1.3pt) node[above=4pt] {$F$};
\fill (E) circle (1.3pt) node[below right=3pt] {$E$};
\fill (B) circle (1.7pt) node[below left=4pt] {$B$};
\fill (O) circle (1.7pt) node[above left=3pt] {$O$};
\fill (A) circle (1.7pt) node[above=4pt] {$A$};
\fill (C) circle (1.7pt) node[right=8pt] {$C$};
\draw[<->,teal!80!black] (0,-2.35)--(1.5,-2.35) node[midway,fill=white,inner sep=3pt] {$x_2$};
\draw[<->] (0,-3.0)--(3,-3.0) node[midway,fill=white,inner sep=3pt] {$x_1$};
\node at (1.5,-3.6) {$\triangle OBF\sim\triangle OCE$};
\node at (1.5,-4.45) {$\dfrac{OB}{OC}=\dfrac{OF}{OE}=\dfrac{x_2}{x_1}$};
\end{scope}
\begin{scope}[shift={(6.2,0)}]
\coordinate (O) at (0,0); \coordinate (A) at (3,1.8);
\coordinate (B) at (1.5,-.9); \coordinate (C) at (3,-1.8);
\coordinate (D) at ({12/17},{36/85}); \coordinate (H) at ({24/17},{72/85});
\fill[orange!12] (O)--(A)--(C)--cycle;
\fill[teal!23] (O)--(A)--(B)--cycle;
\draw[blue!65!black,very thick] (O)--(A);
\draw[violet!75!black,thick,dashed] (O)--(C);
\draw[orange!70!black] (A)--(C);
\draw[teal!80!black,thick] (A)--(B);
\draw[red!65!black,dashed,thick] (B)--(D) (C)--(H);
\pic[draw,black!65,angle radius=3mm] {right angle=B--D--A};
\pic[draw,black!65,angle radius=3mm] {right angle=C--H--A};
\fill (D) circle (1.3pt) node[above left=3pt] {$D$};
\fill (H) circle (1.3pt) node[above left=3pt] {$H$};
\pic[draw,black!65,angle radius=5.5mm] {angle=B--O--A};
\path (B)--(D) node[midway,sloped,above=4pt,text=red!65!black] {$h_B$};
\node[red!65!black] at (2.7,-.45) {$h_C$};
\fill (O) circle (1.7pt) node[above left=3pt] {$O$};
\fill (A) circle (1.7pt) node[above=4pt] {$A$};
\fill (B) circle (1.7pt) node[below left=4pt] {$B$};
\fill (C) circle (1.7pt) node[right=8pt] {$C$};
\node at (1.55,-2.65) {$BD=h_B,\quad CH=h_C$};
\node at (1.55,-3.6) {$\triangle OBD\sim\triangle OCH$};
\node at (1.55,-4.45) {$\dfrac{h_B}{h_C}=\dfrac{BD}{CH}=\dfrac{OB}{OC}$};
\end{scope}

\end{tikzpicture}
\end{center}
先取 $x_1,x_2>0$。图中 $A=(x_1,kx_1)$，$C=(x_1,-kx_1)$，
$B=(x_2,-kx_2)$。左图直接给出
\[
 AC=2kx_1,\qquad S_{OAC}=\frac12\times2kx_1\times x_1=kx_1^2.
\]
\textbf{第一次相似（左图）：}作 $BF\perp x$ 轴、$CE\perp x$ 轴，
则 $OF=x_2,OE=x_1$。直角三角形 $OBF,OCE$ 在 $O$ 处共角，故
\[
 \triangle OBF\sim\triangle OCE
 \quad\Longrightarrow\quad \frac{OB}{OC}=\frac{OF}{OE}=\frac{x_2}{x_1}.
\]
\textbf{第二次相似（右图）：}作 $BD\perp OA,CH\perp OA$，
令 $BD=h_B,CH=h_C$。当 $k\ne1$ 时，直角三角形 $OBD,OCH$ 在 $O$ 处的角相同，故
\[
 \frac{h_B}{h_C}=\frac{OB}{OC}=\frac{x_2}{x_1},\qquad
 S_{OAB}=S_{OAC}\frac{h_B}{h_C}=kx_1x_2.
\]
若 $k=1$，则 $OA\perp OC$，直接有 $h_B=OB,h_C=OC$，比值仍成立。
图取 $k=3/5,x_1=3,x_2=3/2$；当 $B$ 在 $C$ 外侧时，比例大于 $1$，结论同样成立。

\textbf{推广到其他象限：}$S_{OAC}=kx_1^2$ 仍成立，长度比改用
$OB/OC=|x_2/x_1|$；即使 $B,C$ 在反向射线上，面积比仍是正数。
\[
\boxed{S_{OAB}=k|x_1x_2|}\quad\xrightarrow{\ k=b/a\ }\quad
\boxed{S_{AOB}=\frac ba|x_1x_2|}.
\]
一个 $x_1$ 来自水平高；另一个由面积比替换成 $|x_2|$。
\textbf{这一面积关系先属于两条正比例函数，再用于双曲线渐近线。}
\end{explanationbox}

\clearpage
\begin{explanationbox}
\textbf{纵向伸缩：把两条斜线变成直角。}
\begin{center}% C088 PDF teaching figure; exact coordinates, equal x/y units.
\begin{tikzpicture}[x=0.9cm,y=0.9cm,>=stealth,font=\small]
\draw[rounded corners=5pt,black!25] (-1,-4.7) rectangle (12.6,5.6);
\node[font=\bfseries] at (5.8,4.95) {只伸缩纵坐标：两条渐近线变成直角};
\begin{scope}[shift={(0,0)}]
\node at (2,3.9) {原图：$k=2/3$};
\draw[->,black!45] (-0.5,0)--(4.3,0) node[right] {$x$};
\draw[->,black!45] (0,-3.45)--(0,2.6) node[above] {$y$};
\draw[blue!65!black,thick] (0,0)--(3.3,2.1999999999999997);
\draw[violet!75!black,thick,dashed] (0,0)--(3.3,-2.1999999999999997);
\coordinate (A) at (3,2.0); \coordinate (B) at (1.5,-1.0);
\fill[teal!18] (0,0)--(A)--(B)--cycle;
\draw[orange!85!black,thick] (A)--(B);
\draw[dotted,black!45] (A)--(3,0) (B)--(1.5,0);
\fill (0,0) circle (1.6pt) node[below left=4pt] {$O$};
\fill (A) circle (1.8pt) node[right=7pt] {$A$};
\fill (B) circle (1.8pt) node[below left=5pt] {$B$};
\node[below] at (3,0) {$x_1$}; \node[above] at (1.5,0) {$x_2$};
\node at (2,-4) {$S=3$};
\end{scope}
\begin{scope}[shift={(7,0)}]
\node at (2,3.9) {新图：$Y=y/k$，纵向放大 $3/2$};
\draw[->,black!45] (-0.5,0)--(4.3,0) node[right] {$x$};
\draw[->,black!45] (0,-3.45)--(0,2.6) node[above] {$Y$};
\draw[blue!65!black,thick] (0,0)--(3.3,3.3);
\draw[violet!75!black,thick,dashed] (0,0)--(3.3,-3.3);
\coordinate (A) at (3,3); \coordinate (B) at (1.5,-1.5);
\fill[teal!18] (0,0)--(A)--(B)--cycle;
\draw[orange!85!black,thick] (A)--(B);
\draw[dotted,black!45] (A)--(3,0) (B)--(1.5,0);
\fill (0,0) circle (1.6pt) node[below left=4pt] {$O$};
\fill (A) circle (1.8pt) node[right=7pt] {$A'$};
\fill (B) circle (1.8pt) node[below left=5pt] {$B'$};
\node[below] at (3,0) {$x_1$}; \node[above] at (1.5,0) {$x_2$};
\node at (2,-4) {$S'=9/2$，$S=kS'$};
\coordinate (A) at (3,3); \coordinate (B) at (1.5,-1.5); \coordinate (O) at (0,0);
\pic[draw,black!70,angle radius=4mm] {right angle=B--O--A};
\end{scope}

\end{tikzpicture}
\end{center}
保持横坐标不动，只将纵坐标变成 $Y=y/k$，于是
\[
 A'=(x_1,x_1),\qquad B'=(x_2,-x_2).
\]
新图中两条直线 $Y=x$、$Y=-x$ 垂直，不依赖截点所在的象限。
因此
\[
 S'=\frac12\bigl(\sqrt2|x_1|\bigr)\bigl(\sqrt2|x_2|\bigr)=|x_1x_2|.
\]
图中 $k=2/3,x_1=3,x_2=3/2$：纵向放大 $3/2$ 倍，
面积也由 $3$ 放大为 $9/2$；恢复原图要乘 $k$：
\[
\boxed{S=kS'=k|x_1x_2|}.
\]
\textbf{不要直接认定原图的两条渐近线垂直。}
只有 $a=b$、即 $k=1$ 时，原图就是等轴双曲线的直角配置。
一般情况下，直角出现在\emph{伸缩后的图}里。
\end{explanationbox}

\clearpage
\begin{explanationbox}
\textbf{深一层的坐标变换：两渐近线变成坐标轴。}
\begin{center}% C088 PDF teaching figure; exact coordinates, equal x/y units.
\begin{tikzpicture}[x=0.9cm,y=0.9cm,>=stealth,font=\small]
\draw[rounded corners=5pt,black!25] (-1.1,-3.5) rectangle (12.5,5.2);
\node[font=\bfseries] at (5.7,4.55) {渐近线变成坐标轴：$X=x/a+y/b,\ Y=x/a-y/b$};
\node at (1.8,3.5) {原坐标：$a=3,b=2$};
\node at (8.8,3.5) {新坐标：双曲线为 $XY=1$};
\begin{scope}
\draw[->,black!45] (-0.5,0)--(4.35,0) node[right] {$x$};
\draw[->,black!45] (0,-1.5)--(0,2.3) node[above] {$y$};
\draw[blue!65!black,thick] (0,0)--(3.6,2.4);
\draw[violet!75!black,thick,dashed] (0,0)--(3.6,-2.4);
\fill[teal!18] (0,0)--(3,2)--(1.5,-1)--cycle;
\draw[orange!85!black,thick] (3,2)--(1.5,-1);
\draw[black!65,dotted,thick,domain=-1.7:1.9,samples=70,variable=\t] plot ({3*sqrt(1+\t*\t/4)},\t);
\fill (0,0) circle (1.7pt) node[below left=4pt] {$O$};
\fill (3,2) circle (1.7pt) node[above left=5pt] {$A$};
\fill (1.5,-1) circle (1.7pt) node[below left=5pt] {$B$};
\node at (1.8,-2.9) {$S=3$};
\end{scope}
\begin{scope}[shift={(7,0)}]
\draw[->,black!45] (-0.5,0)--(4.35,0) node[right] {$X$};
\draw[->,black!45] (0,-1.5)--(0,2.3) node[above] {$Y$};
\fill[teal!18] (0,0)--(2,0)--(0,1)--cycle;
\draw[orange!85!black,thick] (2,0)--(0,1);
\draw[black!65,dotted,thick,domain=.45:3.7,samples=70,variable=\t] plot (\t,{1/\t});
\fill (0,0) circle (1.7pt) node[below left=4pt] {$O$};
\fill (2,0) circle (1.7pt) node[below=6pt] {$A'$};
\fill (0,1) circle (1.7pt) node[left=6pt] {$B'$};
\node at (1.8,-2.2) {$A'=(2,0),\quad B'=(0,1)$};
\node at (1.8,-2.9) {$S'=1,\quad S=(ab/2)S'$};
\end{scope}

\end{tikzpicture}
\end{center}
设
\[
 X=\frac xa+\frac yb,\qquad Y=\frac xa-\frac yb.
\]
于是 $XY=x^2/a^2-y^2/b^2$。双曲线变成 $XY=1$；
$g_+$ 变成 $Y=0$，$g_-$ 变成 $X=0$。注意图中 $A,B$ 始终是渐近线截点：
\[
 A'=\left(\frac{2x_1}{a},0\right),\qquad
 B'=\left(0,\frac{2x_2}{a}\right),\qquad
 S'=\frac{2|x_1x_2|}{a^2}.
\]
\textbf{面积单位不能混用。}逆变换为
$x=\tfrac a2(X+Y)$、$y=\tfrac b2(X-Y)$，其行列式绝对值为 $ab/2$。
故
\[
 S=\frac{ab}{2}S'=\frac{ab}{2}\cdot\frac{2|x_1x_2|}{a^2}
 =\boxed{\frac ba|x_1x_2|}.
\]
图取 $a=3,b=2,x_1=3,x_2=3/2$，新图面积 $1$，原图面积 $3$；
两图使用同样的数值刻度，阴影面积的差别正好体现缩放因子 $3$。
\end{explanationbox}

\clearpage
\begin{explanationbox}
\textbf{联系反比例函数，也要分清用了哪个三角形。}
\begin{center}% C088 PDF teaching figure; exact coordinates, equal x/y units.
\begin{tikzpicture}[x=0.85cm,y=0.85cm,>=stealth,font=\small]
\draw[rounded corners=5pt,black!25] (-2.6,-2.6) rectangle (11.9,5.1);
\node[font=\bfseries] at (4.65,4.449999999999999) {坐标乘积都能表示面积，所用三角形不同};
\node at (1,3.65) {单个反比例函数点：$xy=6$};
\node at (8,3.65) {两条渐近线上的截点};
\begin{scope}
\draw[->,black!45] (-1.85,0)--(4.2,0) node[right] {$x$};
\draw[->,black!45] (0,-0.8)--(0,2.45) node[above] {$y$};
\fill[orange!12] (0,0) rectangle (3,2);
\fill[teal!23] (0,0)--(3,0)--(3,2)--cycle;
\draw[black!40,dashed] (0,2)--(3,2);
\draw[blue!65!black,thick,domain=2.2:4,samples=60,variable=\t] plot (\t,{6/\t});
\draw[teal!80!black,thick] (0,0)--(3,2)--(3,0)--cycle;
\fill (0,0) circle (1.7pt);\node at (-.35,-.65) {$O$};
\fill (3,2) circle (1.7pt) node[above right=6pt] {$P$};
\fill (3,0) circle (1.7pt) node[below=5pt] {$Q$};
\node at (1,-1.7) {$S_{OPQ}=|xy|/2=|q|/2$};
\end{scope}
\begin{scope}[shift={(7,0)}]
\draw[->,black!45] (-1.85,0)--(4.2,0) node[right] {$x$};
\draw[->,black!45] (0,-0.8)--(0,2.45) node[above] {$y$};
\draw[blue!65!black,thick] (-1.7,-1.1333333333333333)--(3.8,2.533333333333333);
\draw[violet!75!black,thick,dashed] (-1.7,1.1333333333333333)--(1.5,-1);
\fill[teal!23] (0,0)--(3,2)--(-1.5,1)--cycle;
\draw[orange!85!black,thick] (-1.8,.933333)--(3.5,2.111111);
\fill (0,0) circle (1.7pt);\node at (-.35,-.65) {$O$};
\fill (3,2) circle (1.7pt) node[above left=5pt] {$A$};
\fill (-1.5,1) circle (1.7pt) node[above right=5pt] {$B$};
\node at (1,-1.7) {$S_{OAB}=k|x_1x_2|$};
\end{scope}

\end{tikzpicture}
\end{center}
左图中 $P=(3,2)$ 在 $xy=6$ 上，$Q=(3,0)$。
轴平行矩形面积为 $|xy|=6$，阴影三角形是它的一半：
\[
 S_{OPQ}=\frac12|xy|=\frac{|q|}{2}=3.
\]
右图中 $A=(3,2)$、$B=(-3/2,1)$ 分别位于 $y=\pm(2/3)x$，
它们是两条不同直线上的截点，而非同一反比例函数上的两个点。
\[
 S_{OAB}=\frac23\left|3\left(-\frac32\right)\right|=3.
\]
两例的面积数值碰巧相同，但\textbf{左图是一个曲线上点到坐标轴的投影，右图是两个渐近线截点}。
共同的思想是把坐标乘积转化为面积；不能把左图的定值 $|q|/2$ 当成本题的定值。

\textbf{为什么主公式不再乘 $1/2$？}
在本题中，$y_1=kx_1,y_2=-kx_2$，有向面积的两项为
\[
 x_1y_2-x_2y_1=-kx_1x_2-kx_1x_2=-2kx_1x_2.
\]
这两份相同的乘积抵消面积公式外的 $1/2$。证明一的平行四边形图可直接核对这一关系。
\end{explanationbox}
